\section{Introduction}
\label{sec:intro}

Large language models (LLMs) have become the dominant paradigm for knowledge-intensive question answering, yet they remain vulnerable to hallucination and knowledge staleness when answers depend on facts absent from their parameters~\cite{zhao2026survey}. Retrieval-augmented generation (RAG) addresses this by grounding generation in externally retrieved evidence, and the graph-structured variant---GraphRAG---goes a step further by organizing a corpus into a knowledge graph (KG) and reasoning over multi-hop paths between entities~\cite{sun2023thinkongraph,mavromatis2025gnnrag,sun2025searchongraph}. Multi-hop GraphRAG has proven especially valuable for compositional questions whose answers require chaining relations across several hops, a setting where flat vector retrieval fundamentally fails~\cite{pan2023unifying}. The canonical GraphRAG retrieval loop, exemplified by Think-on-Graph (ToG)~\cite{sun2023thinkongraph} and its successors~\cite{ma2024thinkongraph,sun2025searchongraph,liang2025fast}, iteratively expands a frontier of candidate triples, invokes an LLM as a relevance judge to score each candidate, prunes to a fixed beam, and repeats for a predetermined number of hops.

This fixed-budget retrieval policy, however, ignores a fundamental property of real query streams: \emph{questions vary enormously in reasoning difficulty}. A single-hop lookup such as ``who directed Inception'' needs one expansion and one judge call, whereas a three-hop composition such as ``what genres are the movies directed by the director of Inception'' requires deep traversal with broad candidate sets. Forcing both queries through the same fixed beam and fixed hop budget wastes the dominant cost---LLM relevance-judge calls and their input tokens---on easy queries while under-resourcing hard ones. Recent work has begun to address this heterogeneity. \emph{Use-Graph-When-It-Needs}~\cite{dong2026graph} trains a detector to decide whether a query needs the graph, and \emph{GraphRAG-Router}~\cite{fan2026graphragrouter} learns a reinforcement-learned routing policy over GraphRAG backends. These approaches, however, introduce a \emph{training requirement}: they require labeled routing data, a separate training stage, and a learned policy that may not transfer across knowledge graphs or domains. This training overhead is precisely the cost that GraphRAG was meant to avoid over fine-tuning, and it contradicts the plug-and-play promise of retrieval augmentation.

We observe that the signals needed to adapt retrieval \emph{already exist} in the GraphRAG pipeline at no additional cost. Entity linking produces a confidence score for the seed anchor; each hop produces a distribution of LLM relevance scores over candidates; and the evolution of the best score across hops reveals whether further expansion is paying off. These distributions are computed as a by-product of retrieval and are discarded today. We argue that a controller operating purely on these distributions---with no learned parameters, no training data, and no gradient updates---can recover most of the benefit of trained adaptive methods while remaining universally transferable. The key insight is that retrieval difficulty is a property of the \emph{query-pipeline interaction}, not of the query alone, and this interaction is already observable through score distributions.

We propose \method{} (\textbf{A}daptive \textbf{Qu}ery-budgeted \textbf{a}daptive \textbf{RAG}), a training-free, query-adaptive retrieval controller for GraphRAG. \method{} makes three per-query, per-hop decisions using only distributional signals: (i) a \emph{use-graph decision} that routes low-confidence-anchored queries to a cheap graph-free answer path; (ii) an \emph{adaptive beam allocator} that widens the beam when hop-level relevance scores are flat (uncertain) and narrows it when one candidate dominates; and (iii) a \emph{marginal-gain early-stopping rule} with a simple theoretical justification that halts expansion once the best per-hop relevance score stops improving. Crucially, \method{} and the fixed-budget baseline share one retrieval engine; the controller is the \emph{only} variable, so our ablations isolate each component's contribution by construction rather than by experimental confounding.

We evaluate \method{} on MetaQA~\cite{zhang2017variational}, the standard multi-hop KGQA benchmark used by ToG and GNN-RAG, across 1/2/3-hop reasoning depths and two generator scales (1.5B and 7B parameters). Our headline result is that on 2-hop reasoning with a 7B generator---the setting that best balances difficulty and judge reliability---\method{} \emph{strictly Pareto-dominates} the fixed ToG-style baseline: higher F1, higher Hit@1, and lower retrieval cost, simultaneously, with statistical significance. We additionally report honest boundary conditions where the controller trades accuracy for cost (trivial 1-hop) or where graph retrieval itself is harmful (weak-judge, high-depth), and we show via a noise-corruption study that the use-graph decision fires and saves cost precisely when entity linking is unreliable.

\paragraph{Contributions.}
\begin{itemize}[leftmargin=*,nosep]
  \item We identify that GraphRAG applies a uniform expensive retrieval budget to heterogeneous queries, and that the distributions needed to adapt this budget are already computed and discarded by existing pipelines (\Cref{sec:method}).
  \item We propose \method{}, a training-free controller with three components---use-graph decision, adaptive beam, marginal-gain early stopping---each justified by a distributional signal and a simple stopping bound. \method{} requires zero training, zero learned parameters, and runs on a single 32\,GB GPU (\Cref{sec:method}).
  \item We show that \method{} strictly Pareto-dominates the fixed ToG-style baseline at 2-hop with a 7B generator (F1 $+10.3\%$, cost $-23.2\%$, $p{=}0.049$), and improves F1 by $23.5\%$ at $19\%$ lower cost with a 1.5B generator ($p{=}0.004$). Ablations isolate each component's contribution (\Cref{sec:exp}).
  \item We characterize the method's boundaries honestly: trivial queries trade accuracy for cost, weak-judge high-depth settings make graph retrieval harmful, and the use-graph decision fires under noisy linking to save up to $31\%$ cost (\Cref{sec:exp-analysis}).
\end{itemize}
